\documentclass[compacto,10pt]{pubciencias-notas}

% -- Paquetes adicionales
\usepackage{comandosEPN}

% -- Datos del libro
\autor{Departamento de Formación Básica}
\materia{Álgebra Lineal}
\nota{Semana 10 Clase 01}
\fecha{Semestre 2025-A}

\logouno[1.1cm]{Logo_EPN.eps}
\logodos[1.7cm]{Logo_DFB.eps}

\definecolor{colordef}{cmyk}{0.81,0.62,0.00,0.22}

\begin{document}

\encabezado

\vspace*{-4mm}

\setcounter{section}{3}
\section{Espacios Vectoriales}
\setcounter{subsection}{7}
\subsection{Suma de espacios vectoriales}

\begin{defi}[Suma de subespacios]
    Sea $(E,+,\cdot,\K)$ un espacio vectorial y $W_1,W_2\subseteq E$ subespacios vectoriales de $E$. Se define la suma $W_1+W_2$ por:
    \[
        W_1 + W_2 =\left\{v\in E: v=w_1+w_2, \, \text{ con }\, w_1\in W_1, w_2\in W_2\right\}. 
    \]
\end{defi}
\begin{prop}
    Sean $(E,+,\cdot,\K)$ un espacio vectorial y $W_1,W_2\subseteq E$ subespacios vectoriales de $E$. Se tiene que $W_1 + W_2$ es un subespacio vectorial de $E$.
\end{prop}

\begin{proof}[Demostración.]\hspace{0pt} De forma inmediata se tiene que $W_1+W_2$ es diferente del vacío, pues $0_E\in W_1+W_2$, dado que:
\[0_E = 0_E + 0_E\]
donde $0_E\in W_1$ y $0_E\in W_2$ (por ser subespacios vectoriales de E).\newline

Por otro lado, para $\lambda\in\mathbb{K}$, $a\in W_1+W_2$ y $b\in W_1+W_2$ existen $a_1\in W_1$, $a_2\in W_2$, $b_1\in W_1$, $b_2\in W_2$ tales que:
\[
a = a_1 + a_2\quad \text{ y }\quad b = b_1 + b_2
\]
entonces,
\begin{align*}
    \lambda a + b &= \lambda(a_1+a_2) + (b_1+b_2)\\
                  &= (\lambda a_1 + b_1) + (\lambda a_2 + b_2)
\end{align*}    
Lo que muestra la proposición.
\end{proof}


\begin{prop}
    Sean $(E,+,\cdot,\K)$ un espacio vectorial y $W_1,W_2\subseteq E$ dos subespacios vectoriales de $E$. Sean $B_1$ y $B_2$ bases de $W_1$ y $W_2$, respectivamente, se tiene que:
    \[
        W_1 + W_2 = \spn(B_1 \cup B_2).
    \]
\end{prop}

\begin{ejem}\hspace{0pt} Considere el espacio vectorial $\R^{3}$, donde:
\[
    W_{1} = \{(x,0,0)\in \R^3: x \in \R \} 
    \texty
    W_{2} = \{(0,y,0)\in \R^3: y \in \R\}
\]
son subespacios vectoriales, se tiene que: 
\[
    B_{1} = \{(1, 0, 0)\}
    \texty
    B_{2} = \{(0, 1, 0)\}
\]
son bases de $W_{1}$ y $W_{2}$, respectivamente. Entonces:
\begin{align*}
    W_1+W_2
    &=\spn(B_{1} \cup B_{2})\\
    &= \spn(\{(1,0,0), (0,1,0)\})\\
    &= \{(x,y,z)\in\R^{3}\,:\, (x,y,z) = \alpha_{1}(1,0,0) + \alpha_{2}(0,1,0),\, \alpha_{1},\alpha_{2}\in\R \}\\
    &= \{(x,y,z)\in\R^{3}\,:\, (x,y,z) = (\alpha_{1},\alpha_{2},0),\, \alpha_{1},\alpha_{2}\in\R \}\\
    &= \{(x,y,z)\in\R^{3}\,:\, \alpha_{1} = x, \alpha_{2} = y, z=0,\, \alpha_{1},\alpha_{2}\in\R \}\\
    &= \{(x,y,z)\in \R^3:  z = 0\}\\
    &= \{(x,y,0)\in \R^3:  x,y\in \R \}
\end{align*}
\end{ejem}

\begin{ejer}\hspace{0pt} Considere el espacio vectorial $\mathbb{R}^{2\times 2}$ y los subespacios vectoriales:
\[S=\left\{\begin{pmatrix} a & b\\ c & d \end{pmatrix} : a+2c=0\right\}\]
\[T=\left\{\begin{pmatrix} a & b\\ c & d \end{pmatrix} : d-a=0\right\}\]
Halle $S+T$ y su dimensión.
\end{ejer}

\begin{teo}[Teorema de Grassmann]
    Sean $(E,+,\cdot,\K)$ un espacio vectorial de dimensión finita y $W_1,W_2\subseteq E$ dos subespacios vectoriales de $E$. Se tiene que
    \[
        \dim(W_1 + W_2) = \dim(W_1) + \dim(W_2) - \dim(W_1 \cap W_2).
    \]
\end{teo}

\begin{obs}
De acuerdo con el teorema anterior concluimos que en un espacio vectorial $E$ de dimensión finita, cualquier subespacio $W$ verifica que
\[
\dim(W) \leq \dim(E).
\]
\end{obs}

\begin{ejem}\hspace{0pt} Considere los subespacios vectoriales:
\[
    W_{1} = \{(x,0,0)\in \R^3: x \in \R \} 
    \texty
    W_{2} = \{(0,y,0)\in \R^3: y \in \R\}.
\]

Se tiene que $\dim(E) = \dim(\R^{3}) = 3$ y que $\dim(W_{1}) = \dim(W_{2}) = 1$ con $\dim(W_{1} + W_{2}) = 2$, por el teorema anterior, concluimos que $\dim(W_{1} \cap W_{2}) = 0$, lo que implica que $W_{1} \cap W_{2} = \{0_{\R^{3}}\}$.
\end{ejem}

\begin{ejer}\hspace{0pt} Para el Ejercicio 1, obtenga la $\dim(S\cap T)$.
\end{ejer}

\begin{defi}[Suma directa]
    Sean $(E,+,\cdot,\K)$ un espacio vectorial y $W_1, W_2\subseteq E$ subespacios vectoriales de $E$. La suma directa de $W_1$ y $W_2$, denotada por $W_1\oplus W_2$ se define por:
    \[
        W_1 \oplus W_2 =\left\{v\in E : (\exists !w_1\in W_1)(\exists !w_2\in W_2)\, \, v=w_1+w_2 \right\} 
    \]
\end{defi}

\begin{advertencia}
Se dice que un espacio vectorial $E$ está en suma directa de dos subespacios $W_1$ y $W_2$ si cualquier elemento de $E$ se puede escribir de manera única como la suma de elementos de $W_1$ y $W_2$.
\end{advertencia}


\begin{teo}
    Sean $(E,+,\cdot,\K)$ un espacio vectorial y  $W_1,W_2\subseteq E$ subespacios vectoriales de $E$, se tiene que:
    \[
        E = W_1 \oplus W_2
    \]
    si y sólo si:
    \[
        E = W_1 + W_2 \texty
        W_1 \cap W_2 = \{0_{E}\}.
    \]
\end{teo}
\begin{ejem}\hspace{0pt}\label{EjemDirecta1} En el espacio vectorial $\R^{3}$ considere los subespacios vectoriales:
\[
    W_{1} = \{(x,y,z)\in \R^3: x+y+z=0\}
    \texty 
    W_{2} = \{(x,y,z)\in \R^3: x=z=0\}.
\]
Se tiene que:
\begin{align*}
    W_{1} &= \{(x,y,-x-y): x,y\in\R\}\\
    &= \{x(1,0,-1)+y(0,1,-1): x,y\in\R\}\\
    &= \langle \{\underbrace{(1,0,-1), (0,1,-1)}_{:=B_1}\}\rangle
\end{align*}
con $B_1$ un conjunto linealmente independiente, es decir, $B_1=\{(1,0,-1), (0,1,-1)\}$ es una base de $W_1$. Además,
\begin{align*}
    W_{2} &= \{(0,y,0): y\in\R\}\\
    &= \{y(0,1,0): y\in\R\}\\
    &= \langle \{\underbrace{(0,1,0)}_{:=B_2}\}\rangle
\end{align*}
con $B_2$ un conjunto linealmente independiente, es decir, $B_2=\{(0,1,0)\}$ es una base de $W_2$.
Por tanto:
\begin{align*}
    W_1+W_2 &= \spn(B_1\cup B_2)\\
    &= \{(x,y,z)\in\R^3 : (x,y,z)=\lambda_1 (1,0,-1) + \lambda_2 (0,1,-1) + \lambda_3 (0,1,0) \}\\
    &= \{(x,y,z)\in\R^3 : (\lambda_1, \lambda_2+\lambda_3, -\lambda_1-\lambda_2)=(x,y,z)\}\\
    &= \R^{3}
\end{align*}
además,
\begin{align*}
    W_1 \cap W_2 &= \{(x,y,z)\in\R^3 : x+y+z=0, x=z=0\}\\
    &= \{(x,y,z)\in\R^3 : x=y=z=0\}\\
    &= (0,0,0)
\end{align*}
de lo anterior se concluye:
\[
    \R^{3} = W_{1} \oplus W_{2}.
\]
\end{ejem}

\begin{ejer}\hspace{0pt} En el espacio vectorial $\mathbb{R}^{2\times 2}$ considere los subespacios vectoriales:
\[W_1 = \left\{A\in \mathbb{R}^{2\times 2} : A=A^T\right\}\]
\[W_2 = \left\{A\in \mathbb{R}^{2\times 2} : A=-A^T\right\}\]
Verifique que: $\mathbb{R}^{2\times 2} = W_1 \oplus W_2$.
\end{ejer}

\end{document}